% Generated by parameter_sensitivity_report.py
% This analysis evaluates a separable 1-D Taylor model applied independently along x and y. It is not the joint full 2-D Taylor model.
\section{Parameter sensitivity analysis}
\label{sec:parameter-sensitivity}

\paragraph{Scope.}
This analysis evaluates a separable 1-D Taylor model applied independently along $x$ and $y$. It is not the joint full 2-D Taylor model. Results use parameter-matched clean references and seed 42. They describe matched-clean noise stability for the tested image and must not be interpreted as absolute gradient accuracy or a population-level estimate.

All descriptor metrics use fixed Robust HOG with L2-Hys normalization and Huber weighting (delta=2.0, window=5). Runtime definition: Median of 3 noisy gradient-plus-HOG runs; per-configuration design-matrix setup and clean-reference computation excluded. Residual definition: Unweighted RMS Taylor-fit residual over both per-axis systems on the common analysis ROI; diagnostic, not correntropy loss. Residual and iteration count are diagnostics and are not components of the composite score.

\paragraph{Effect of the per-axis sample count.}
With $h=1$, increasing $n$ from 16 to 32 reduced correntropy RMSE by 60.3\% under 5\% salt-and-pepper noise and by 65.3\% under 10\% noise. The corresponding runtime increases were 54.2\% and 42.1\%. Thus, $n=32$ maximized robustness among the tested values, whereas $n=16$ provided a more economical correntropy configuration.

\paragraph{Effect of sample spacing.}
With $n=8$, increasing $h$ improved all matched-clean stability metrics. The final step from $h=3$ to $h=4$ reduced correntropy RMSE by 21.6\% at 5\% noise and 24.0\% at 10\% noise, but the cosine gains were only 0.0049 and 0.0052. Accordingly, $h=4$ is the best tested stability setting, while $h=3$ is a more conservative localization and conditioning compromise.

\paragraph{Solver trade-off.}
Direct and iterative squared loss produced nearly identical HOG stability, while the iterative method required several times more runtime. For the larger-support tested configurations, correntropy typically produced substantially lower RMSE and HOG error and higher cosine similarity under impulsive noise, but its runtime was tens of times larger than direct squared loss. It is therefore justified when robustness is prioritized, not as an unconditional replacement for the direct solver.

\paragraph{Recommended settings.}
The maximum-robustness tested setting is correntropy with $n=32,h=1$; a cost-aware alternative is $n=16,h=1$. In the separate spacing sweep $h=4$ was best among the tested values. The experiment is one-factor-at-a-time, so it does not support claiming $(n,h)=(32,4)$ as a joint optimum. The requested scenario-balanced composite also ranked correntropy loss with $n=32,h=1$ first (score 0.8661), but this score is candidate-set and weight dependent.
